\begin{tikzpicture}[node distance=1.7cm and 2.6cm,
startstop/.style={rectangle, rounded corners, draw, thick,
minimum width=4.0cm, minimum height=1.0cm,
font=\small, align=center, fill=red!20},
process/.style={rectangle, rounded corners, draw, thick,
minimum width=5.0cm, minimum height=1.2cm,
text width=4.7cm, align=center,
font=\small, fill=blue!15},
smallprocess/.style={rectangle, rounded corners, draw, thick,
minimum width=4.3cm, minimum height=1.2cm,
text width=4.0cm, align=center,
font=\small, fill=blue!15},
io/.style={trapezium, trapezium left angle=70, trapezium right angle=110,
draw, thick, minimum width=4.4cm, minimum height=1.2cm,
align=center, font=\small, fill=green!20},
decision/.style={diamond, draw, thick, aspect=2,
minimum width=3.4cm, minimum height=1.8cm,
align=center, font=\small, fill=yellow!20},
subprocess/.style={rectangle, draw, thick, double,
minimum width=5.0cm, minimum height=1.2cm,
text width=4.7cm, align=center,
font=\small, fill=purple!20},
arrow/.style={->, thick, >=stealth},
]

% Title
\node[font=\Large\bfseries] at (0,14.3) {System Operation Flowchart (Multi-Target Mode)};

% Row 1
\node[startstop] (start) at (0,12.5) {System Start};

% Row 2
\node[process] (init) [below=1.5cm of start] {Initialize ESP-NOW Protocol};
\node[process] (config) [below=1.4cm of init] {Load Configuration (Time Slots, Node IDs)};
\node[process] (sync) [below=1.4cm of config] {Synchronize All Nodes to a Common Time Base};

% Row 3
\node[decision] (running) [below=2.2cm of sync] {System Running?};
\node[startstop] (stop) [left=4cm of running] {Stop System};

% Row 4
\node[process] (slotcheck) [below=2.2cm of running] {Check Current Time Slot};

% Row 5
\node[decision] (slotactive) [below=2.0cm of slotcheck]{Is This Node's Time Slot?};

% Target branch
\node[process] (targetTx) [below left=3cm and 4.4cm of slotactive]{Target Node:\\Broadcast ESP-NOW Packet};
\node[process] (targetWait) [below=3.4cm of targetTx]{Wait for Next Assigned Slot (100 ms)};

% Anchor branch
\node[process] (anchorRx) [below right=3cm and 4.4cm of slotactive]{Anchor Node:\\Receive the Broadcast Packet};
\node[smallprocess] (anchorRSSI) [below=1.2cm of anchorRx]{Measure RSSI and Timestamp};
\node[smallprocess] (anchorSend) [below=1.2cm of anchorRSSI]{Send Processed Data to the Gateway};

% Gateway block
\node[subprocess] (gateway) [below=3.5cm of slotactive]{Gateway: Aggregate All Anchor Data};
\node[decision] (allresp) [below=2cm of gateway]{All Anchors Responded?};
\node[smallprocess] (waitTimeout) [right=4cm of allresp, yshift=-1.5cm]{Wait for Timeout Window};

% Estimation & output
\node[subprocess] (estimate) [below=2cm of allresp]{Estimate Position (Based on Aggregated RSSI)};
\node[io] (output) [below=2cm of estimate]{Output Estimated Position to PC / Logger};
\node[process] (store) [below=1.7cm of output]{Store Position Data in Database};

% Arrows
\draw[arrow] (start) -- (init);
\draw[arrow] (init) -- (config);
\draw[arrow] (config) -- (sync);
\draw[arrow] (sync) -- (running);
\draw[arrow] (running) -- node[midway, above] {No} (stop);
\draw[arrow] (running) -- node[midway, right] {Yes} (slotcheck);
\draw[arrow] (slotcheck) -- (slotactive);

% Target branch routing (clean, no overlaps)
\draw[arrow] (slotactive.west) -- ++(-1,0)|- node[pos=0.2, above] {Yes (Target)} (targetTx);
\draw[arrow] (targetTx) -- (targetWait);
\draw[arrow] (targetWait.west) -- ++(-2.5,0)|- (slotcheck.west);

% Anchor branch routing
\draw[arrow] (slotactive.east) -- ++(1,0)|- node[pos=0.25, above] {Yes (Anchor)} (anchorRx);
\draw[arrow] (anchorRx) -- (anchorRSSI);
\draw[arrow] (anchorRSSI) -- (anchorSend);
\draw[arrow] (anchorSend.east) -- ++(2.4,0)|- (gateway.east);

% Gateway → Decision
\draw[arrow] (gateway) -- (allresp);

% ===== FIXED SECTION (No through-node lines) =====
\draw[<->, thick, >=stealth] (allresp.east) -- (waitTimeout.west);

% Yes-branch
\draw[arrow] (allresp.south) -- (estimate);

% Downstream
\draw[arrow] (estimate) -- (output);
\draw[arrow] (output) -- (store);

% Loop back cleanly
\draw[arrow] (store.west) -- ++(-3.3,0)|- (slotcheck.west);

\end{tikzpicture}
